\documentclass[10pt,twocolumn]{article}
\usepackage[a4paper,margin=1.9cm,columnsep=0.7cm]{geometry}
\usepackage{mathptmx}
\usepackage[T1]{fontenc}
\usepackage{microtype}
\usepackage{graphicx}
\usepackage{booktabs}
\usepackage{caption}
\usepackage{authblk}
\usepackage[hidelinks]{hyperref}
\usepackage{xcolor}

\captionsetup{font=small,labelfont=bf,skip=5pt}
\raggedbottom
\setlength{\parindent}{1.1em}
\definecolor{mono}{HTML}{2F6F9F}
\definecolor{dia}{HTML}{C25A3C}
\definecolor{m3gray}{HTML}{8A8A8A}

\title{\vspace{-1.2em}\bfseries Monologue or Dialogue?\\[2pt]
\large Measuring Conversational Structure in the Reasoning Traces of Large Language Models}
\author[ ]{Oskar Engl}
\affil[ ]{\small Johannes Kepler University Linz, Austria}
\affil[ ]{\small \textit{oskar.engl@ik.me}}
\date{\small Working version, July 2026}

\begin{document}
\twocolumn[
\begin{@twocolumnfalse}
\maketitle
\begin{abstract}
\noindent
A reasoning model such as DeepSeek-R1 produces a long ``thinking'' trace before it gives an answer.
That trace can be read as one mind working step by step, or as several voices exchanging proposals
and objections. This paper asks which reading the text itself supports once it is measured. Reasoning
models are recent, so there is little prior work on the voice structure of their traces, and no trace
comes labelled with its true structure; we therefore compare against a small grid of corpora whose
structure is known: human and AI, monologue and dialogue. We test three text-analysis methods and
carry forward those that correctly classify all four known corpora as monologue or dialogue. Two of
the three do so; the third, a coarse document-level author count, fails to recognise multiple voices
even in a genuine four-way debate, so it does not judge the traces. Both carried-forward methods place
the traces on the dialogue side of the scale, most clearly in how the writing style shifts from
sentence to sentence and in how often a question is followed by an answer. We are explicit about how
much weight each method can bear, and flag that one of the two depends on an automatic labeller of
only moderate reliability.
\medskip
\end{abstract}
\end{@twocolumnfalse}
]

\begin{figure*}[t]
\centering
\includegraphics[width=0.96\textwidth]{fin_fig_grid_v2.pdf}
\caption{\textbf{Study design.} Four corpora of known voice structure (left) provide a validity
check: we carry forward the methods that correctly classify all four known cells as monologue or
dialogue. Two of the three methods do so; the coarse document-level author count does not, and is set
aside. Carried-forward methods are applied to the DeepSeek-R1 traces (\texttt{ai-reason}) and place
them on a monologue--dialogue scale anchored by the two same-task AI cells; the human cells serve only
to check that a method works.}
\label{fig:grid}
\end{figure*}

\section{Introduction}
Reasoning-oriented training has become standard in frontier models. Given a problem, a model such as
DeepSeek-R1 \cite{guo2025} first produces a long stretch of intermediate text, often called a
reasoning trace \cite{kambhampati2025}, and only then gives its final answer. These traces have a
notable feature: each is written by a single model in one pass, yet a trace often contains questions
and their answers, objections and concessions, and sudden reversals such as ``Wait, that is not
right.'' In writing by people, features like these can mark a real exchange between participants,
though they can also mark a single inner voice. This raises the question: is a reasoning trace's
textual shape closer to that of a monologue or of a dialogue?

Kim et al.~\cite{kim2026} argue that reasoning-trained models improve in part because they simulate a
``society of thought,'' a set of internal perspectives that propose, criticise, and reconcile. There
are interesting parallels in older ideas from cognitive science, which have variously understood
reasoning as a social and argumentative activity \cite{mercier2017} and the mind as a set of
interacting agents \cite{minsky1986}. Our study works from an assumption: if a trace is produced by
several voices taking turns, that should leave marks in the writing. Authorship analysis rests on the
observation that writers differ in small, largely involuntary habits, most robustly in how they use
function words, and that these habits are stable within an author and vary between authors
\cite{stamatatos2009}. We therefore treat the question in operational terms: a monologue should read
as the work of one author and a dialogue as the work of several authors, and we ask which of the two
the traces are closer to.

\section{Study design}
Reasoning models are new enough that there is little prior work on the voice structure of their
traces, and no collection of traces comes with its structure labelled. Our approach is to compare
against a grid of corpora whose structure \emph{is} known (Figure~\ref{fig:grid}). We assemble four
such corpora and vary two properties across them: who produced the text (a human or an AI) and whether
it is a monologue or a dialogue. The four cells are human monologue (think-aloud protocols, in which a
person speaks their thoughts aloud while solving a task \cite{wurgaft2025}), human dialogue (group
deliberation from DeliData \cite{karadzhov2023}), AI monologue (a single Qwen-2.5-7B model solving a
puzzle on its own), and AI dialogue (four separate copies of Qwen-2.5-7B, each given a role, taking
turns in a debate). The fifth cell is DeepSeek-R1 itself: the model is given the same puzzle, and we
take its native thinking trace --- the reasoning it produces before its final answer --- as the
document. This is the cell whose structure we want to determine.

The grid functions as a validity check. We define the check at the level of the cell and apply it
identically to every method: does the method place each of the four known cells on the correct side of
the monologue--dialogue axis? A method is carried forward to the traces only if it classifies all four
known cells correctly. The rationale runs both ways: a method that cannot recognise multiple voices in
a genuine four-way debate is in no position to rule them out in a reasoning trace, and a method that
reports several voices in one person's think-aloud protocol is in no position to affirm them in a
trace either.

We hold the task fixed where possible. The AI monologue, AI dialogue, and DeepSeek-R1 cells all solve
the same puzzle, Game of 24 \cite{yao2023}, a small arithmetic puzzle in which four given numbers must
be combined with the four basic operations to reach 24, and the human monologue cell uses the same
puzzle. This rules out a method separating the cells by their subject matter rather than by their
voice structure. The human dialogue cell is the one exception, since it uses the Wason card task,
another compact logic puzzle, because no group-deliberation corpus exists for Game of 24; its role is
only to test whether a method can detect human dialogue at all, so it never anchors the placement. We
also removed speaker names, role titles, and letter labels from both dialogue corpora before analysis,
so that a method cannot detect dialogue by identifying who is speaking instead of reading the
structure of the exchange.

\section{Methods}
The three methods read deliberately different signals, so that agreement among them carries more
weight than any single method alone. Each turns its signal into a monologue-or-dialogue call for a
cell; Section~4 reports which methods pass the validity check and what they say about the traces.

\textbf{Method 1: counting authors across the document} \cite{nath2019}. This method treats a
document as a whole. It divides the text into paragraph-sized windows, describes each window by how
often the 50 most common words occur in it, and groups similar windows together. The number of groups
is read as the number of authors, which we reduce to a binary decision: one author (monologue) or more
than one (dialogue). A cell counts as correctly classified if its documents fall on the right side of
that decision on balance.

\textbf{Method 2: detecting style changes between sentences} \cite{schmidt2025}. This method operates
sentence by sentence. For each pair of neighbouring sentences it estimates the probability that the
two were written by different authors, using an encoder trained to capture writing style while
ignoring topic \cite{patel2025}. Longer documents contain more sentence pairs, which affects the raw
score, so we adjust for length before comparing cells. Each document receives a single number, which
we call its \emph{lean}, indicating whether its style-change pattern lies closer to the dialogue cells
or to the monologue cells; a cell is classified by the direction of its average lean.

\textbf{Method 3: measuring question-and-answer structure} \cite{duran2022,schegloff1973}. This
method reads the text as a sequence of conversational moves. A language model tags each utterance with
its function using CAMS, a standard scheme whose labels include a question, an answer, an agreement,
or a suggestion. From these tags we measure how often an \emph{adjacency pair} is completed. An
adjacency pair is a two-part exchange in which the first part makes the second one due, such as a
question followed by an answer, or a proposal followed by its acceptance. A single writer has little
reason to produce completed pairs, so their presence is evidence of genuine dialogue; a cell is
classified by whether its dialogue documents complete more pairs than its monologue documents. Because
Method~3's labels are produced by a language model acting as a tagger, their reliability is a separate
concern, which we measure against a human annotator and return to in Sections~4 and~6.

\section{Results}
\textbf{The validity check.} Of the three methods, two --- Method~2 and Method~3 --- correctly classify
all four known cells as monologue or dialogue, and are carried forward to the traces. Method~1 does
not: it misclassifies both dialogue cells, and is set aside. We take each method in turn.

\textbf{Method 1 fails the check.} Read as a summary number, its balanced accuracy is $0.6088$ (95\%
confidence interval $0.596$ to $0.621$), above the chance level of $0.50$; but that single number
hides a one-sided failure. On single-author text the method is nearly perfect, correctly labelling
$0.975$ of the human monologue documents and $0.998$ of the AI monologue documents. On multi-author
text it collapses: it recovers only $0.423$ of the AI debates --- so a majority of genuine four-way
debates are read as a single author --- and \emph{none} of the human debates, calling all $500$ of
them single-author. In cell terms, it places both known dialogue cells on the monologue side. A method
that cannot see four speakers in a four-speaker debate, and that never once flags the human debate as
multi-voiced, has not shown it can detect voices at all, so it does not judge the traces.

\textbf{A note on the deviation from pre-registration.} We originally pre-registered a separate
statistical gate for each method rather than the single cell-level check used here. Under those
per-method gates, Method~1's balanced accuracy cleared a chance baseline and would have ``passed,''
while Method~3 was blocked by a label-reliability threshold ($\kappa \ge 0.70$) discussed below. On
inspection those gates behaved pathologically: they admitted a method that labels genuine four-way
debate as single-author, and excluded a method that separates all four known cells correctly. The
per-method gates were only ever an operationalisation of the idea stated from the outset --- that a
method must separate monologue from dialogue on the known corpora before it is trusted on the traces
--- and that operationalisation did not match the idea. We therefore report against the simpler,
uniform criterion we set out to use: does the method correctly classify all four known cells? This is
a stated deviation, not a silent one; the pre-registered gates and the reason for departing from them
are recorded here so a reader can weigh both.

\textbf{Method 2 passes the check cleanly.} All four cells fall on the correct side: the monologue
cells lean negative and the dialogue cells positive (average lean $-0.0105$ for the monologue cells
against $0.0125$ for the dialogue cells), and the separation is strong, with Cliff's $\delta=0.6211$
(Cliff's $\delta$ is a rank-based measure of how cleanly two groups separate, from $0$ at complete
overlap to $1$ at no overlap). Applied to the traces, Method~2 places them further toward dialogue
than either AI anchor.

\textbf{Method 3 passes the check, with a caveat about its labels.} On the classification the check
asks for, it performs well. Its completion rate separates dialogue from monologue in the expected
direction (Cliff's $\delta=0.4717$): within each producer, the dialogue cell completes more pairs than
the monologue cell --- $0.235$ against $0.096$ for the human cells, and $0.025$ against $0.003$ for the
AI cells. Applied to the traces, its completion rate of $0.035$ lies closer to the AI dialogue anchor
($0.025$) than to the AI monologue anchor ($0.003$), a dialogue-side placement.

The caveat is that this classification rests on labels from an automatic tagger, and separating the
cells is not the same as labelling every utterance correctly: a tagger could be systematically wrong
in a consistent way and still separate the cells. We measured the tagger against a human annotator
using Cohen's $\kappa$, an agreement score that corrects for chance matches, and obtained
$\kappa=0.54$ --- below the $0.70$ we had pre-registered. We do not treat this as disqualifying, for a
reason we can state but not fully settle here. CAMS uses $28$ fine-grained labels, many of them close
relatives such as \texttt{inform} and \texttt{suggest} that are easily confused even by careful human
coders, and unweighted $\kappa$ charges a near miss between two similar tags as heavily as a
disagreement between two unrelated ones. The CAMS authors make the same point and recommend a measure
that gives partial credit for near misses; under such a measure agreement rises to a tree-weighted
$0.62$, and the raw share of labels on which tagger and human agree is $0.74$. These figures argue
that the labels are \emph{reasonable}, not that they are validated: they remain below our threshold,
and we do not claim the labels are error-free. The honest resolution --- an adjudicated
double-annotation --- is left to future work (Section~6). We therefore carry Method~3 forward on the
strength of its cell-level separation, while flagging that its label reliability is the weakest link
in its chain.

\begin{figure*}[t]
\centering
\includegraphics[width=0.90\textwidth]{fin_fig_split_v2.pdf}
\caption{\textbf{Placement of the traces by each carried-forward method} on the monologue--dialogue
scale, anchored so \texttt{ai-mono}$=0$ and \texttt{ai-dia}$=1$. Both methods place the DeepSeek-R1
traces past the dialogue pole and at nearly the same point ($1.45$ and $1.47$); Method~3 is daggered
to mark that its placement depends on an automatic labeller of only moderate reliability (Section~4).
Method~1 is not shown here: it failed the validity check and does not place the traces (Table~1).}
\label{fig:split}
\end{figure*}

\textbf{Across the methods} (Table~\ref{tab:verdict} and Figure~\ref{fig:split}), the two methods that
pass the validity check agree: both place the traces on the dialogue side. Method~2, reading style
changes between sentences, places them firmly past the dialogue anchor; Method~3, reading completed
question-and-answer pairs, places them past it too. The one method that pointed toward monologue,
Method~1, is not a dissenting vote we weigh against these two --- it is a method that failed the check,
and its reading of the traces is reported only to show it was not quietly dropped. For what it is
worth, even Method~1 places the traces in ambiguous territory rather than cleanly at the monologue
pole, but its verdict carries no weight, because it cannot classify the known dialogue cells to begin
with.

\begin{table*}[t]
\centering
\small
\setlength{\tabcolsep}{8pt}
\caption{Each method against the single validity check --- does it place all four known cells on the
correct side of the monologue/dialogue axis? --- and, for the two that pass, where it places the
traces. Method~3 passes the check but depends on an automatic labeller of moderate reliability
($\dagger$; Section~4).}
\label{tab:verdict}
\begin{tabular}{@{}lccl@{}}
\toprule
Method & Classifies all four known cells? & Carried forward? & Verdict on traces \\
\midrule
1: author count & no & no & --- \\
2: style switch & yes & yes & \textcolor{dia}{dialogue} \\
3: adjacency pairs & yes & yes$^{\dagger}$ & \textcolor{dia}{dialogue}$^{\dagger}$ \\
\bottomrule
\end{tabular}
\end{table*}

\section{Discussion}
The claim here concerns the surface of the text, that is, the patterns a reader or a tool can observe
in the writing. It says nothing directly about what occurs inside the model. On that surface, the
picture leans toward dialogue. The traces change writing style from sentence to sentence in the manner
of dialogue, and they complete question-and-answer pairs at a rate above the dialogue anchor.

The scope of this claim is limited. A dialogue-leaning reading does not show that a model contains
several minds, just as a novel full of dialogue does not have several authors. It shows that the trace
is organised like a conversation, which is a measurable fact in itself, and a natural surface
counterpart to the internal multi-agent behaviour that Kim et al.~\cite{kim2026} describe from the
inside, through interpretability. The outside view and the inside view point the same way, which is
reassuring, though the present evidence stands on its own as a statement about surface structure.

The strength of the finding survives even a sceptical reading of Method~3. Taken together, the two
validated methods agree that the traces are dialogue-like, and that is the case we make. But suppose a
reader sets Method~3 aside entirely on the grounds that its labels are only moderately reliable. The
validated evidence then reduces to Method~2 --- the method with the cleanest validation --- which
still places the traces on the dialogue side, and \emph{no} validated method places them on the
monologue side. (Method~1 does, but it failed the validity check and is not part of the validated
evidence.) So the dialogue lean does not hinge on trusting the tagger: it holds with Method~3 and
survives without it.

Two caveats temper the dialogue placements, and both concern how a distance was measured rather than
its direction. Method~2's distance to the monologue side is partly an extrapolation, because the
traces are much longer than the monologue anchor, so it is more accurate to say that they track the
dialogue anchor closely over the range where the two overlap. Method~3's verdict rests on its
completion signal, since its coarser tag-frequency profiles cannot separate the two AI anchors.
Neither caveat changes a direction; each limits only how far past the pole the placement should be
read.

\section{Limitations}
Several limitations should be noted. The target is a single reasoning model, DeepSeek-R1, on a single
family of puzzles, so the verdict concerns this setting and does not yet generalise to reasoning
traces at large. The two anchors were generated with a different base model (Qwen-2.5-7B) than the
target, so part of the measured distance may reflect the change of model as well as voice structure,
and the placement is best read as ``closer to which pole'' than as an exact distance. Method~2's
monologue comparison rests on extrapolation, as noted above.

Method~3's labels come from an automatic tagger whose agreement with a human annotator ($\kappa=0.54$
unweighted; tree-weighted $0.62$; raw agreement $0.74$) is only moderate. The validity check we apply
confirms that the tagger \emph{separates} the known cells in the right direction, but not that every
label is correct, and these two are distinct; a systematic labelling bias could separate the cells
without the labels being trustworthy in detail. The rigorous remedy is an adjudicated
double-annotation, with two independent coders and a resolution step, which we leave to future work.
Until then, Method~3 should be read as corroborating evidence rather than as an independently
validated labelling.

Finally, a note on method more than on data. We report against a single cell-level validity check
applied identically to all three methods, having departed from a more elaborate set of per-method
gates fixed in advance (Section~4). The reason for the departure is stated rather than hidden, and the
uniform check is the sharper instrument: a method earns the right to judge the traces exactly when it
can tell monologue from dialogue on data whose structure we already know.

\begin{thebibliography}{99}
\small
\setlength{\itemsep}{0.6pt}\setlength{\parskip}{0pt}
\bibitem{guo2025} D.\ Guo et al. DeepSeek-R1 incentivizes reasoning in LLMs through reinforcement
learning. \textit{Nature} 645:633--638, 2025.
\bibitem{kambhampati2025} S.\ Kambhampati et al. Position: Stop anthropomorphizing intermediate
tokens as reasoning/thinking traces! \textit{arXiv:2504.09762}, 2025.
\bibitem{kim2026} J.\ Kim, S.\ Lai, N.\ Scherrer, B.\ Ag\"uera y Arcas, J.\ Evans. Reasoning models
generate societies of thought. \textit{arXiv:2601.10825}, 2026.
\bibitem{mercier2017} H.\ Mercier, D.\ Sperber. \textit{The Enigma of Reason.} Harvard Univ.\
Press, 2017.
\bibitem{minsky1986} M.\ Minsky. \textit{The Society of Mind.} Simon \& Schuster, 1986.
\bibitem{stamatatos2009} E.\ Stamatatos. A survey of modern authorship attribution methods.
\textit{JASIST} 60(3):538--556, 2009.
\bibitem{wurgaft2025} D.\ Wurgaft et al. Scaling up the think-aloud method. \textit{Proc.\ CogSci},
2025.
\bibitem{karadzhov2023} G.\ Karadzhov, T.\ Stafford, A.\ Vlachos. DeliData: a dataset for
deliberation in multi-party problem solving. \textit{Proc.\ ACM HCI} 7(CSCW2), 2023.
\bibitem{yao2023} S.\ Yao et al. Tree of thoughts: deliberate problem solving with LLMs.
\textit{NeurIPS} 36, 2023.
\bibitem{nath2019} S.\ Nath. Style change detection by threshold-based and window-merge clustering.
\textit{CLEF Working Notes}, 2019.
\bibitem{schmidt2025} G.\ Schmidt et al. Style change detection using a sequential sentence-pair
classifier. \textit{CLEF Working Notes}, 2025. \textit{arXiv:2508.00675}.
\bibitem{patel2025} A.\ Patel et al. StyleDistance: content-independent style embeddings with
synthetic parallel examples. \textit{NAACL}, 2025.
\bibitem{duran2022} N.\ Duran, S.\ Battle, J.\ Smith. Inter-annotator agreement using the
Conversation Analysis Modelling Schema. \textit{Comm.\ Methods \& Measures} 16(3):182--214, 2022.
\bibitem{schegloff1973} E.\ Schegloff, H.\ Sacks. Opening up closings. \textit{Semiotica}
8(4):289--327, 1973.
\end{thebibliography}

\end{document}
